% Modern editorial, markup and figure/layout contributions: CC-BY-SA-4.0.
% Historical text and illustrations: Leonhard Euler (1744), public domain.
% Component scope and upstream attribution: see RIGHTS.md.
% 校录范围：原印刷页 17--23（PDF 页 20--26）。
\sourcepage{17}
atque ita porro quæcunque formula debeat in curva quæsita esse maxima vel minima, ea semper erit hujus formæ $\int Z\,dx$, scilicet integrale quantitatis finitæ cujusdam $Z$ in differentiale $dx$ ductæ. Debet autem $Z$ ejusmodi esse quantitas, ut si æquatio statuatur inter $x$ \& $y$, integrale $\int Z\,dx$ determinatum obtineat valorem: ex quo $Z$ erit functio quantitatum $x$, $y$, \& inde pendentium $p$, $q$, $r$, \&c. vel algebraica sive determinata, vel præterea ipsa in se complectetur formulas integrales indeterminatas; quod discrimen probe est tenendum. Ita si maximi minimive formula $W$ fuerit $\int y\,dx$, vel $\int y\,dx\sqrt{(1+pp)}$; quantitas $Z$ erit algebraica, at si sit $W=\int yx\,dx\int y\,dx$, tum erit $Z=yx\int y\,dx$, hoc est ipsa quantitas $Z$ erit indeterminata, cujus valor nisi relatio inter $x$ \& $y$ detur, exhiberi nequit. Quin etiam evenire potest, ut valor ipsius $Z$ hujusmodi formula evoluta exprimi nequeat, sed tantum per æquationem differentialem demum erui debeat, ut si fuerit $dZ=y\,dx+ZZ\,dx$; ex qua æquatione valor ipsius $Z$ per $x$ \& $y$ nequidem exhiberi potest. Hinc igitur tria nascuntur genera formularum $\int Z\,dx$, quæ in curvis quæsitis maxima vel minima fieri debent. Quorum primum eas complectitur formulas, in quibus $Z$ est functio algebraica seu determinata ipsarum $x$, $y$, \& $p$, $q$, $r$, \&c. Ad secundum genus referimus eas formulas, in quibus quantitas $Z$ ipsa insuper formulas integrales involvit. In tertio autem genere continentur eæ formulæ, in quibus valor ipsius $Z$ per æquationem differentialem, cujus integratio non constat, determinatur.

\subsection*{Propositio II. Theorema.}
\originalsection{38}\textit{Si fuerit $amz$ curva, in qua valor formulæ $\int Z\,dx$ sit maximus vel minimus, atque $Z$ sit functio algebraica seu determinata ipsarum $x$, $y$, $p$, $q$, $r$, \&c. tum ejusdem curvæ quæcunque portio $mn$ eadem gaudebit prærogativa, ut pro ea ad suam abscissam $MN$ relata, valor ipsius $\int Z\,dx$ sit pariter maximus vel minimus.}

\subsection*{Demonstratio.}
Valor\sourcepage{18} formulæ $\int Z\,dx$ pro abscissa $AZ$ est aggregatum omnium valorum ejusdem formulæ, qui singulis abscissæ $AZ$ portionibus respondent. Quod si ergo abscissa $AZ$ in partes quotcunque, quarum una sit $MN$, divisa concipiatur, atque ad singulas partes hasce valor formulæ $\int Z\,dx$ exhibeatur; summa omnium horum valorum præbebit valorem formulæ $\int Z\,dx$, qui toti abscissæ $AZ$ convenit; \& qui erit maximus vel minimus. Quoniam autem $Z$ ponitur functio algebraica ipsarum $x$, $y$, $p$, $q$, \&c. valor formulæ $\int Z\,dx$ respondens abscissæ portioni $MN$, a sola portionis curvæ respondentis $mn$ indole pendebit, idemque manebit, utcunque reliquæ partes $am$ \& $nz$ varientur; singularum enim litterarum $x$, $y$, $p$, $q$, \&c. valores per solam curvæ portionem $mn$ determinantur. Si ergo formulæ $\int Z\,dx$ valores, qui conveniunt abscissæ portionibus $AM$, $MN$, $NZ$, ponantur $P$, $Q$ \& $R$, quantitates hæ $P$, $Q$, \& $R$, a se mutuo non pendebunt. Quare cum earum aggregatum $P+Q+R$ sit maximum vel minimum, etiam unaquæque maximi minimive proprietate prædita sit necesse est. Hanc ob rem, si in curva $amz$ formula $\int Z\,dx$ maximum minimumve habeat valorem, \& quantitas $Z$ sit functio algebraica ipsarum $x$, $y$, $p$, $q$, \&c. tum etiam, pro qualibet illius curvæ portione, eadem formula $\int Z\,dx$ maximi minimive proprietate gaudebit. Q. E. D.

\subsection*{Coroll. I.}
\originalsection{39}Quod si ergo curva fuerit inventa $amz$, quæ pro abscissa data $AZ$, habeat valorem formulæ $\int Z\,dx$ maximum vel minimum, atque $Z$ sit functio algebraica seu determinata, tum etiam ejusdem curvæ qualibet portio, respectu abscissæ suæ respondentis, eadem maximi minimive proprietate gaudebit.

\subsection*{Coroll. II.}
\originalsection{40}\sourcepage{19}In hujusmodi igitur Problematibus, ubi tale maximum minimumve quæritur, non opus est quantitatem abscissæ, cui maximum minimumve respondeat, definire; sed si, pro una quacunque abscissa, formula $\int Z\,dx$ sit maximum vel minimum, tum eadem pro quacunque alia abscissa eadem proprietate gaudebit.

\subsection*{Coroll. III.}
\originalsection{41}Hujusmodi igitur Problemata resolventur, si singulæ curvæ quæsitæ particulæ ita determinentur, ut pro iis valor formulæ $\int Z\,dx$ fiat maximus vel minimus. Tum enim simul tota curva, \& quæcunque ejus portio, pariter eadem maximi minimive proprietate erit instructa.

\subsection*{Scholion.}
\originalsection{42}Proprietas hæc, qua gaudent curvæ in quibus istiusmodi formulæ $\int Z\,dx$, ubi $Z$ est functio algebraica seu determinata ipsarum $x$, $y$, $p$, $q$, \&c. sunt maximum vel minimum, est maximi momenti; ea enim innititur universa methodus hujus generis Problemata resolvendi. Ideo autem potissimum hanc Propositionem afferre visum est, ne ea proprietas, quæ his tantum formulis $\int Z\,dx$, ubi $Z$ est functio vel algebraica vel determinata, est propria, omnium omnino formularum quæ proponi possunt communis esse putetur: in sequente enim Propositione demonstrabimus, si in $Z$ insint formulæ integrales; tum eandem proprietatem non amplius locum habere: ex quo simul natura hujusmodi quæstionum clarius intelligetur. Hujus autem præsentis Propositionis demonstratio ex eo petita est fundamento, quod valor formulæ $\int Z\,dx$, siquidem $Z$ est functio vel algebraica vel determinata ipsarum $x$, $y$, $p$, $q$, $r$, \&c. qui convenit cuicunque abscissæ portioni $MN$, a sola curvæ portione respondente $mn$ pendeat, neque a reliqua curva, vel anteriore $am$, vel posteriore $nz$ afficiatur: quæ ratio cessat, si in $Z$ insint \sourcepage{20}formulæ integrales indeterminatæ. Valores enim quantitatum $x$, $y$, $p$, $q$, $r$, \&c. qui pro arcu curvæ $mn$ obtinent, tantum a positione elementorum hujus arcus $mn$, atque elementis aliquot contiguis, quæ arcum finitæ quantitatis non constituunt, pendent; ex quo etiam quantitas ex iis litteris utcunque composita per solam arcus $mn$ indolem determinabitur, nisi adfuerint quantitates integrales, cujusmodi sunt $\int y\,dx$, quæ totam aream anteriorem $AamM$ introduceret, vel $\int dx\sqrt{(1+pp)}$, quæ totum arcum præcedentem $am$ involveret. Hinc igitur distinctius intelligitur, quid per functionem determinatam ipsarum $x$, $y$, $p$, $q$, $r$, \&c. denotare velimus: Functio scilicet determinata ita est comparata, ut, pro quovis loco, a præsentibus valoribus litterarum $x$, $y$, $p$, $q$, \&c. tantum pendeat, neque valores earum anteriores in se complectatur. Functio autem indeterminata est talis, cujus valor in quovis loco, non ex solis valoribus quos hæ litteræ $x$, $y$, $p$, $q$, \&c. in isto loco obtinent determinari potest, sed insuper omnes valores ad sui determinationem requirit, quos istæ litteræ in omnibus locis anterioribus obtinuerunt. Ita patet, omnes functiones algebraicas esse simul determinatas; præterea vero etiam omnes functiones transcendentes, quæ a relatione inter $x$ \& $y$ non pendent, sunt determinatæ, cujusmodi sunt, $l\sqrt{(xx+yy)}$, $e^{p y}$, $A\sin.\frac{p y}{q}$; quarum valores in quovis loco ex valoribus litterarum, quos in hoc solo loco obtinent, assignari possunt. Quando autem in functione quapiam insunt formulæ integrales indeterminatæ, quæ a mutua relatione inter $x$ \& $y$, quam ubique tenent, pendent, tum earum valor, in dato loco, non ex valoribus, quos hæ litteræ in isto loco habent, cognosci potest, sed insuper omnes valores in locis quibusque anterioribus nosse oportet, hoc est generalem relationem inter coordinatas $x$ \& $y$: talesque functiones vocamus indeterminatas; quippe quæ toto cœlo diversæ sunt ab iis, quas determinatas appellavimus.

\subsection*{Propositio III. Theorema.}
\originalsection{43}\sourcepage{21}\textit{Si fuerit $amz$ curva abscissæ $AZ$ respondens, in qua $\int Z\,dx$ sit maximum vel minimum; in $Z$ autem contineantur formulæ integrales indeterminatæ; tum eadem maximi minimive proprietas non cadit in quamlibet curvæ portionem, sed toti tantum curvæ abscissæ $AZ$ respondenti propria erit.}

\subsection*{Demonstratio.} Concipiatur tota curva $amz$, pro qua $\int Z\,dx$ est maximum vel minimum, in duas partes quasvis divisa per applicatam $Mm$; sitque formulæ $\int Z\,dx$ valor conveniens portioni $am$ $=P$, ejusdem autem formulæ valor pro altera portione $mz$ sit $=Q$: pro tota igitur curva $amz$ valor formulæ $\int Z\,dx$ erit $=P+Q$, quem ponimus esse maximum vel minimum. Quo autem omnem ambiguitatem tollamus, totamque rem distinctius proponere queamus; ponamus $P+Q$ esse maximum: quod enim de maximo demonstrabitur, idem de minimo facile intelligetur. Quod si jam valor ipsius $Q$ a valore ipsius $P$ non dependeret, tum aggregatum $P+Q$ maximum esse non posset, nisi simul uterque valor $P$ \& $Q$ seorsim sit maximus. At nostro casu, quo quantitas $Z$ in se continet formulas integrales indeterminatas, valor ipsius $Q$ non tantum a curvæ portione $mz$ ad quam refertur pendebit, sed simul a tota curva anteriore $am$; atque adeo a valore ipsius $P$. Nunc dicimus, ad id ut $P+Q$ sit maximum, non requiri, ut valor ipsius $P$ sit maximus. Ponamus enim portionem curvæ $am$ ita esse comparatam, ut pro ea $P$ sit maximum, \& aliquantillum mutari concipiatur portio curvæ $am$, ita ut valor formulæ $\int Z\,dx$ minor evadat, puta $=P-p$: fieri utique poterit ut ex hac mutatione valor ipsius $Q$ crescat, quod incrementum ponatur $q$: eritque, mutata aliquantillum portione $am$, ita ut pro ea $\int Z\,dx$ non amplius sit maximum, valor formulæ $\int Z\,dx$ pro tota curva $amz =P-p+Q+q$. Cum igitur evenire queat ut sit $q>p$, intelligitur \sourcepage{22}formulam $\int Z\,dx$ pro tota curva $amz$ maximam esse posse, etiamsi maxima non sit pro qualibet portione $am$. Q. E. D.

\subsection*{Coroll. I.}
\originalsection{44}Quando ergo curva fuerit inventa, quæ, pro data abscissa $AZ$, habeat valorem formulæ $\int Z\,dx$ maximum vel minimum, \& $Z$ sit functio indeterminata; tum non sequitur quamlibet curvæ inventæ portionem eadem maximi minimive proprietate fore præditam.

\subsection*{Coroll. II.}
\originalsection{45}In resolutione igitur hujusmodi Problematum, in quibus curva quæritur, quæ pro data abscissa $AZ$ habeat $\int Z\,dx$ maximum vel minimum, perpetuo ad totius abscissæ propositæ quantitatem erit respiciendum, atque maximum vel minimum ad eam tantum, non vero ad ejus quamlibet portionem, accommodari debebit.

\subsection*{Coroll. III.}
\originalsection{46}Maximum igitur hinc patet discrimen, quod inter formulas $\int Z\,dx$, in quibus $Z$ functio est determinata vel indeterminata, intercedit; simulque autem Methodorum diversitas intelligitur, quibus ad resolutiones quæstionum, in quibus hujusmodi formularum maximi minimive valores requiruntur, uti oportebit.

\subsection*{Scholion.}
\originalsection{47}Ex demonstratione hujus Propositionis non quidem necessario sequitur, si pro data abscissa $AZ$ curva habeat formulam $\int Z\,dx$ maximam vel minimam, tum singulas ejus portiones eadem hac prærogativa gaudere; verumtamen satis intelligitur, quoties eadem proprietas in singulas portiones competat, id \sourcepage{23}casu evenire.
Hincque nihilominus summe necessarium est, solutionem perpetuo ad totam propositam abscissam accommodare. Interim tamen, in Problematibus ad methodum relativam pertinentibus evenire potest, ut formulas $\int Z\,dx$, in quibus $Z$ sit functio indeterminata, quasi determinata esset tractare liceat. Hoc scilicet accidit, si inter omnes tantum curvas in quibus formulæ illæ integrales indeterminatæ quæ in $Z$ insunt æquales obtinent valores, ea desideretur, in qua $\int Z\,dx$ sit maximum vel minimum: hoc enim casu formulæ illæ integrales indeterminatæ fieri censendæ sunt determinatæ. Ita si, inter omnes curvas ejusdem longitudinis, determinanda sit ea in qua sit $\int Z\,dx$ maximum vel minimum, atque in $Z$, præter quantitates determinatas, insit arcus curvæ $\int dx\sqrt{(1+pp)}$; hic, quia in omnibus curvis ex quibus quæsitam definire oportet, eundem obtinet valorem, instar functionis determinatæ tractari poterit: Hæc autem cuncta in sequentibus clarius explicabuntur.

\Needspace{14\baselineskip}
\subsection*{Hypothesis II.}
\originalsection{48}\marginnote{\footnotesize\hyperlink{fig:2}{Fig. 2.}}\textit{Si curvæ abscissa $AZ$ in elementa innumerabilia infinite parva \& inter se æqualia dissecetur, cujusmodi sunt $IK$, $KL$, $LM$, \&c. atque portio quæcunque $AM$ vocetur $x$, cui respondeat functio quæcunque variabilis $F$, eandem functionem $F$, quatenus refertur ad puncta abscissæ vel sequentia $N$, $O$, $P$, $Q$, \&c. vel antecedentia $L$, $K$, $I$, \&c. ita denotabimus, ut sit valor ipsius functionis, qui pro puncto $M$ est $=F$, ut sequitur.
\[
\left.\begin{array}{rcl}
\text{pro }N&=&F'\\
\text{pro }O&=&F''\\
\text{pro }P&=&F'''\\
\text{pro }Q&=&F^{\mathrm{iv}}\\
\text{pro }R&=&F^{\mathrm{v}}\\
&\text{\&c.}
\end{array}\right\}
\qquad\text{pro punctis abscissæ sequentibus}
\]}

